\documentclass[11pt]{article}
\usepackage[utf8]{inputenc}
\usepackage{amsmath, amssymb, amsthm, bm}
\usepackage{geometry}
\geometry{margin=1in}

\newtheorem{definition}{Definition}
\newtheorem{theorem}{Theorem}
\newtheorem{proposition}{Proposition}
\newtheorem{corollary}{Corollary}

\title{\textbf{Why Common Stock Prevents State-Space Pricing:\\Rank Deficiency, Waterfall Double-Counting, and the Constructive Path to Arrow--Debreu Completeness}}
\author{Javaad Ali}
\date{}

\begin{document}

\maketitle

\begin{abstract}
The Fundamental Theorem of Asset Pricing ($p_t = \mathbb{E}_t[M_{t+1} x_{t+1}]$) is exact when traded securities span the underlying state space $\Omega$. We argue that the central failure of empirical asset pricing is not the theory of state-space valuation, but the reliance on \emph{common stock} as the primary traded primitive. Common stock is a bundled, entity-level residual claim that prevents market completeness through two structural defects: (i) \emph{rank deficiency}, wherein $N$ corporate shells cannot span a $D$-dimensional coupled operational state space ($N \ll D$); and (ii) \emph{legal non-partitioning}, wherein standalone residual claims double-count interdependent receivables across corporate boundaries, generating phantom equity that vanishes under simultaneous default. We construct an explicit mechanism to restore completeness: replacing corporate equity wrappers with \emph{formula-defined, pre-waterfall state primitives} settled via deterministic operational oracles, thereby partitioning aggregate physical cash flows and rendering Arrow--Debreu state prices $q(\omega)$ uniquely invertible.
\end{abstract}

\section{State-Space Completeness and the Equity Bottleneck}

Let $\Omega$ denote the set of future states of the physical and contractual economy, with $\dim(\mathcal{H}_{\mathrm{global}}) = D = \prod_{i=1}^N d_i$ representing the coupled state space across $N$ operating enterprises. By the Fundamental Theorem of Asset Pricing (Cochrane, 2005), if the set of traded payoff vectors $\mathcal{X} = \mathrm{span}\{x_1, \dots, x_K\}$ satisfies $\dim(\mathcal{X}) = D$, there exists a unique stochastic discount factor $M_{t+1} > 0$ and a unique positive state-price vector $\bm{q}_t \in \mathbb{R}_{++}^D$ such that every claim satisfies:
\begin{equation}
    p_{k,t} = \mathbb{E}_t[M_{t+1} x_{k,t+1}] = \sum_{\omega \in \Omega} q_t(\omega) x_{k,t+1}(\omega).
\end{equation}

Modern capital markets, however, do not trade basis vectors of $\Omega$. They trade \emph{common stock}—a legal wrapper that compresses an entire enterprise's internal operations, bilateral contracts, and priority waterfalls into a single scalar residual claim.

\begin{proposition}[Rank Deficiency and Non-Invertibility of Common Equity]
Let $X_{\mathrm{eq}} \in \mathbb{R}^{N \times D}$ be the payoff matrix of $N$ common stocks across $D$ coupled states. Because $N \ll D$ and each row $X_{\mathrm{eq}, i}$ aggregates nonlinear operational states via a lossy contraction $p_i(t) = \mathrm{Re}(\mathrm{Tr}_{-i}[\hat{O}_i \rho(t)])_{g_t}$, we have $\mathrm{rank}(X_{\mathrm{eq}}) \le N \ll D$. Consequently, $\bm{q}_t$ lies in a high-dimensional nullspace ($\dim(\ker(X_{\mathrm{eq}})) \ge D - N$) and cannot be inverted from scalar equity prices.
\end{proposition}

\section{Legal Non-Partitioning: Phantom Equity Under Simultaneous Default}

Consider a network of $N$ corporations linked by bilateral obligations $L_{ij}(\omega)$ (trade payables, OTC derivatives, or sole-source supply dependencies, such as an annuity underwriter $i$ whose workforce constitutes the sole customer base of food vendor $j$). Under standalone accounting, the apparent residual equity payoff of firm $i$ in state $\omega$ assumes counterparty solvency:
\begin{equation}
    E_i^{\mathrm{standalone}}(\omega) = \max\Bigg(0, \; A_i(\omega) + \sum_{j \neq i} L_{ji}^{\mathrm{book}}(\omega) - \sum_{k \neq i} L_{ik}(\omega)\Bigg),
\end{equation}
where $A_i(\omega)$ denotes physical assets and $L_{ji}^{\mathrm{book}}(\omega)$ denotes gross contractual receivables.

\begin{theorem}[Phantom Equity Under Simultaneous Liquidation]
Let $\omega_{\mathrm{sys}} \in \Omega$ denote a systemic clearing state in which all $N$ corporate entities undergo simultaneous bankruptcy liquidation. In $\omega_{\mathrm{sys}}$, contractual receivables collapse to realized recovery rates $\tilde{L}_{ji}(\omega_{\mathrm{sys}}) \le L_{ji}^{\mathrm{book}}(\omega_{\mathrm{sys}})$ determined simultaneously via Eisenberg--Noe fixed-point clearing. Summing standalone equity valuations double-counts interdependent residual profits:
\begin{equation}
    \sum_{i=1}^N E_i^{\mathrm{standalone}}(\omega_{\mathrm{sys}}) - \sum_{i=1}^N E_i^{\mathrm{cleared}}(\omega_{\mathrm{sys}}) = \Delta_{\mathrm{phantom}}(\omega_{\mathrm{sys}}) > 0.
\end{equation}
Common stock violates the Arrow--Debreu resource-feasibility condition: claims do not partition aggregate physical output, because interdependent receivables are counted simultaneously as residual equity across multiple balance sheets.
\end{theorem}

\section{Constructive Path to Completeness: Formula-Defined State Primitives}

To transition from incomplete, double-counted equity shells to exact Arrow--Debreu state-space pricing, traded claims must be legally and mathematically restructured along three constructive principles:

\subsection{1. Pre-Waterfall Physical Partitioning (Eliminating Double-Counting)}
Rather than issuing residual claims at the bottom of circular entity waterfalls, instruments must attach directly to verifiable, non-overlapping physical or net-cleared cash-flow streams prior to inter-company balance-sheet netting. Let $\mathcal{Y}(\omega) = \sum_{i=1}^N A_i(\omega)$ denote aggregate realized physical surplus in state $\omega$. A set of $K$ state primitives $\{x_1, \dots, x_K\}$ forms a \emph{valid legal partition} if and only if:
\begin{equation}
    \sum_{k=1}^K w_k x_k(\omega) \equiv \mathcal{Y}(\omega) \quad \forall \omega \in \Omega,
\end{equation}
with zero unbacked gross receivables ($L_{ji}^{\mathrm{book}}$ netted at clearing), guaranteeing $\Delta_{\mathrm{phantom}}(\omega) \equiv 0$ even under simultaneous default.

\subsection{2. Orthogonal Regime Decomposition (Spanning Asymmetric Operators)}
Recall the asymmetric coupling operator between an institutional anchor $i$ and a physical supplier $j$:
\begin{equation}
    \hat{H}_{ij}\big(x_i(\omega)\big) = \alpha \, \Theta\big(\tau_i - x_i(\omega)\big) \hat{D}_j + \beta \, \sigma\Big(\gamma \big(x_i(\omega) - \mu_i\big)\Big) \hat{I}_j.
\end{equation}
Common stock fuses the linear regime, the downside jump discontinuity $\Theta(\tau_i - x_i)$, and the upside saturation plateau $\sigma(\cdot)$ into a single scalar payoff. Completeness is restored by decomposing $\hat{H}_{ij}$ into \textbf{formula-defined contingent claims} parameterized directly by the observable state variables $\bm{s}(\omega)$:
\begin{align}
    x_{\mathrm{base}}(\omega) &= \min\big(x_i(\omega), \tau_i\big) && \text{(Senior physical throughput claim)}, \\
    x_{\mathrm{sat}}(\omega)  &= \beta \, \sigma\Big(\gamma\big(x_i(\omega) - \mu_i\big)\Big) && \text{(Bounded capacity-ceiling claim)}, \\
    x_{\mathrm{tail}}(\omega) &= \alpha \, \Theta\big(\tau_i - x_i(\omega)\big) && \text{(Binary solvency-breach trigger)}, \\
    x_{\mathrm{net}}(\omega)  &= \mathcal{Y}(\omega) - \big(x_{\mathrm{base}}(\omega) + x_{\mathrm{sat}}(\omega) + x_{\mathrm{tail}}(\omega)\big) && \text{(Cleared residual systemic tranche)}.
\end{align}

\subsection{3. Deterministic Oracle Settlement (Decoupling Fundamentals from $g_t$)}
Instead of quoting scalar spot prices deformed by a time-varying liquidity metric $g_t$, each primitive is defined structurally by its exact payoff formula $x_k = f_k(\bm{s}(\omega))$ verified via deterministic operational data feeds (e.g., audited unit throughput, net cleared swap obligations, or ledgered inventory states).

\begin{corollary}[Exact State-Price Identification]
When the traded security set consists of $D$ linearly independent, pre-waterfall formula primitives $\{f_1(\bm{s}), \dots, f_D(\bm{s})\}$ spanning $\Omega$, the payoff matrix $X_{\mathrm{prim}} \in \mathbb{R}^{D \times D}$ satisfies $\mathrm{rank}(X_{\mathrm{prim}}) = D$ and $\Delta_{\mathrm{phantom}}(\omega) = 0$. Given market clearing prices $\bm{p}_t$, Cochrane's Arrow--Debreu state prices are uniquely and exactingly identified:
\begin{equation}
    \bm{q}_t = X_{\mathrm{prim}}^{-1} \bm{p}_t, \qquad M_{t+1}(\omega) = \frac{q_t(\omega)}{\pi_t(\omega)},
\end{equation}
where $\pi_t(\omega)$ denotes physical state probabilities.
\end{corollary}

\end{document}
